\newcommand{\nomeEstado}{BAHIA}
\newcommand{\siglaEstado}{BA}
\newcommand{\nomeConcessionaria}{NEOENERGIA BAHIA}
\newcommand{\cnpjConcessionaria}{15.139.629/0001-94}
\newcommand{\termoCIP}{Sexto}
\newcommand{\presidenteConcessionaria}{Sr. Thiago Guth.}

\newcommand{\empresaResponsavel}{ABEL}
\newcommand{\nomeMunicipio}{Muniz Ferreira}
\newcommand{\nomeMunicipioMaisculo}{\MakeUppercase{\nomeMunicipio}}
\newcommand{\IPestimada}{3404512}
\newcommand{\cnpjMunicipio}{13.796.461/0001-64}
\newcommand{\sedePrefeitura}{Praça 30 de Julho, nº 168, Centro}
\newcommand{\nomePrefeito}{Gileno Pereira dos Santos}
\newcommand{\generoPrefeitura}{pelo Prefeito}
\newcommand{\numeroContrato}{226/2024}
\newcommand{\quantidadeAditivos}{1}
\newcommand{\legitimidade}{1}
\newcommand{\publicacao}{Portal Nacional de Contratações Públicas}
\newcommand{\dataPublicacao}{05 de agosto de 2025}
\newcommand{\tituloClausula}{Primeiro Termo Aditivo ao Contrato Nº 226/2024}
\newcommand{\clausula}{Primeiro Termo Aditivo ao Contrato Nº 226/2024:

De um lado, o MUNICÍPIO DE MUNIZ FERREIRA, pessoa jurídica de direito público interno, inscrita no CNPJ 13.796.461/0001-64, com sede administrativa Prefeitura Municipal na Praça ACM Junior, no 168, Centro – CEP 44.575-000 Muniz Ferreira – Bahia, representado pelo Prefeito Municipal, Sr. Gileno Pereira dos Santos, brasileiro, maior, casado, CPF 597.084.505-10, RG 05.672.898-00 SSP-BA, doravante designado CONTRATANTE, e do outro a empresa ABEL CUNHA – SOCIEDADE INDIVIDUAL DE ADVOCACIA, Inscrito No CNPJ 29.574.422/0001-52 Localizado na Q CNB 6, no 403, Taguatinga Norte (Taguatinga), Lote 12, Edifício Dona Elvira, CEP: 72.115-065, Brasília- DF, neste ato representado pela Sr. Abel Gomes Cunha, CPF: 991.114.111-04, doravante denominada CONTRATADA, firmam o presente Termo Aditivo ao contrato no226/2024, cujo objeto é a contratação de empresa para prestação de serviços de assessoria e consultoria técnica, jurídico-administrativa especializada na gestão, elaboração de auditorias e laudos técnicos, mediante a conferência das faturas de energia elétrica da administração direta e indireta do Município, elaboração de memorial de cálculo de consumo e potência do parque de iluminação pública, a apuração do modelo tarifário aplicado em cada unidade consumidora, assim como verificação de possíveis isenções indevidas e/ou não repasse da Contribuição de Iluminação Pública (CIP) e/ou não recolhimento do ISS dos prestadores de serviços do setor elétrico, visando a repetição de indébitos decorrentes de cobranças indevidas (a maior) nas contas de energia elétrica de titularidade municipal. Deverá a empresa efetuar levantamento dos créditos a que faz jus o Município, referentes aos pagamentos indevidos à concessionária de energia elétrica, conforme os prazos estabelecidos na Resolução Normativa da ANEEL no 1.000, de 7 de dezembro de 2021, especialmente o art. 323, §2o, inciso I, e suas alterações posteriores, podendo, para tanto, representar administrativamente o Município de Muniz Ferreira - BA, perante a Agência Nacional de Energia Elétrica – ANEEL, perante a distribuidora de energia elétrica e, caso necessário, perante a agência reguladora conveniada do Estado, em todos os processos, reclamações, recursos e demais atos decorrentes de suas competências, com base no art. 124, I, b, da Lei no 14.133/2021, de acordo com o processo licitatório Inexigibilidade 056/2024. (grifou-se)}
\newcommand{\telefone}{(75) 3632-1330}
\newcommand{\email}{falecom@saj.ba.gov.br}
